
\section{Introduction}

\subsection{Motivation}

Standard softmax attention (Vaswani et al., 2017) computes, for each of $H$ 
heads and each pair of positions $(i, j)$, a score:

\[
s_{ij}^{(h)} = \frac{(Q^{(h)} x_i)^\top (K^{(h)} x_j)}{\sqrt{d_h}}
\]

producing an $L \times L$ score matrix per head. This is $\mathcal{O}(L^2 d_h)$ 
in time and $\mathcal{O}(L^2)$ in memory -- the dominant cost for long sequences.

Yet the attention distribution is \textbf{highly concentrated}. Across numerous 
studies (Clark et al., 2019; Voita et al., 2019; Michel et al., 2019), the 
following has been repeatedly observed:

\begin{itemize}
\item Most attention weight is concentrated on a few keys (attention sinks,
  heavy hitters).
\item Many heads are redundant and can be pruned entirely.
\item The sparsity pattern varies dramatically across heads.

\end{itemize}
\textbf{Static sparsity} -- fixed windows (Longformer), fixed strides (BigBird), 
fixed random patterns -- cannot adapt to this variation. A head that needs 
global context is forced into a local window; a head that only needs local 
context wastes capacity on distant keys.

\textbf{Dynamic sparsity} -- top-$k$ selection per query -- adapts to content, but 
existing methods suffer from two problems:

\begin{enumerate}
\item \textbf{Non-differentiability.} Hard top-$k$ selection has zero gradient
  almost everywhere, making it unsuitable for end-to-end training.
\item \textbf{Uniform budget.} Most methods use the same $k$ for all heads, ignoring
  the fact that head importance varies.

\end{enumerate}
\subsection{Problem Statement}

Given attention score matrix $S \in \mathbb{R}^{L \times L}$ per head, we 
seek a sparse approximation $\tilde{S}$ such that:

\[
\tilde{S} \approx S, \qquad
\text{supp}(\tilde{S}) \ll L^2
\]

We impose the following constraints:

\begin{enumerate}
\item \textbf{Differentiability:} End-to-end trainable w.r.t. the sparsity budget.
\item \textbf{Per-head adaptation:} Each head has its own sparsity budget.
\item \textbf{Quality:} $\Delta\text{Acc} \leq 5\%$ on SST-2.
\item \textbf{Memory bounded:} Peak VRAM must not exceed baseline.
\item \textbf{Compatibility:} Works with pretrained weights without retraining.

\end{enumerate}
\subsection{Our Approach: Adaptive Sparse Attention}

We propose \textbf{Adaptive Sparse Attention} (ASA), which addresses the two 
problems above simultaneously:

\begin{enumerate}
\item \textbf{Learnable per-head budget.} Each head maintains a scalar logit
  $\ell_h \in \mathbb{R}$. Its retention budget is:
  \[
k\_h = \lfloor (1 - \sigma(\ell\_h)) \cdot L \rfloor
\]
  where $\sigma$ is the sigmoid. This makes the budget a **continuous,
  differentiable** function of a single parameter.

\item \textbf{Soft top-$k$ masking.} Hard top-$k$ selection is replaced by a soft
  mask:
  \[
\tilde{s}\_{ij} = s\_{ij} + \log \sigma\!\left(\frac{s\_{ij} - \tau\_h}{T}\right)
\]
  where $\tau_h$ is the $k_h$-th largest score for head $h$ and $T = 0.1$
  is a temperature. As $T \to 0$, this converges to hard top-$k$; for finite
  $T$, it is smooth and differentiable.

\item \textbf{Chunked computation.} To bound memory, the query dimension is
  processed in chunks of 64. This keeps peak VRAM at $\mathcal{O}(B H \cdot
  \text{chunk} \cdot L)$ rather than $\mathcal{O}(B H L^2)$.

\end{enumerate}
\subsection{Contributions}

\begin{itemize}
\item \textbf{A learnable per-head sparsity mechanism.} ASA's per-head logits are the
  first differentiable mechanism for learning head-specific sparsity budgets
  without auxiliary losses.
\item \textbf{A soft top-$k$ threshold.} We replace hard top-$k$ with a sigmoid-based
  soft mask that concentrates mass on the retained entries while remaining
  differentiable.
\item \textbf{A production-hardened implementation.} The implementation supports flat,
  nested, and causal architectures with RoPE, GQA, T5-style relative position
  bias, and safe injection.
\item \textbf{A systematic evaluation.} Benchmarked on DistilBERT, RoBERTa, and GPT-2
  on SST-2 with ablations over sparsity ratio, temperature, and head budget
  distribution.
\item \textbf{Head budget analysis.} We show that learned budgets are highly non-uniform,
  with a consistent bimodal distribution across layers.

\end{itemize}
\subsection{Paper Organization}

Section 2 reviews related work. Section 3 provides background on attention and 
sparsity. Section 4 presents ASA. Section 5 describes integration with three 
architectures. Section 6 details the experimental setup. Section 7 reports 
results. Section 8 analyzes head heterogeneity and sparsity patterns. Section 
9 discusses practical guidance. Section 10 lists limitations. Section 11 
concludes.



\subsection{Notation}

\begin{lstlisting}

---
\end{lstlisting}